Django es un marco web de alto nivel para Python, originado en 2005 y
mantenido por la Django Software Foundation \cite{django}. Aporta de serie
un mapeador objeto-relacional (\acs{orm}), un sistema de migraciones de
esquema, gestión de usuarios y permisos y soporte para internacionalización
lo que reduce significativamente la cantidad de código que el desarrollador debe
escribir y mantener. Sobre Django, \ac{drf} añade los componentes
específicos para construir APIs REST: serializadores con validación,
vistas basadas en clases, autenticación intercambiable, paginación y la
generación automática de esquemas OpenAPI \cite{drf}. La elección de
Django y \acs{drf} frente a alternativas como Flask, FastAPI o frameworks
de otros lenguajes obedece a tres factores prácticos: la madurez del
ecosistema y la disponibilidad de extensiones para tareas auxiliares
(generación de diagramas de modelos, panel de administración, soporte
multi-idioma), la familiaridad previa con el marco, y la
convención clara que Django impone sobre la estructura del proyecto,
particularmente útil en un proyecto modular con trece aplicaciones como
el que aquí se aborda.
